\begin{abstract}
Dynamic parameter blending in modular deep learning offers a flexible alternative to static weight averaging, but standard parametric routing networks suffer from a severe \emph{Overfitting-Optimizer Paradox} under low-data calibration regimes (e.g., $N=64$ samples) and exhibit \emph{heterogeneity collapse} under multi-task stream evaluation. We propose \textbf{Gaussian Process Dynamic Routing (GP-DR)}, a mathematically rigorous, non-parametric Bayesian framework that completely bypasses gradient-based optimization of routing parameters. By treating frozen calibration samples as spatial landmarks, GP-DR models dynamic merging coefficients as a closed-form posterior mean under a Gaussian Process prior with a Radial Basis Function (RBF) kernel. While the closed-form posterior predictive variance theoretically offers a metric for epistemic uncertainty, we show that it suffers from a critical unit-sphere variance collapse under realistic noise and that simpler distance heuristics empirically outperform it under representational overlap. We explicitly analyze these mathematical limitations with full scientific transparency. When paired with stream-level \textbf{Micro-Batch Homogenization (MBH)} to prevent cross-task representation interference, GP-DR achieves exceptional performance. On a challenging $14$-layer, $192$-dimensional Isolating Coordinate Sandbox, GP-DR achieves a $72.40\%$ Joint Mean accuracy with zero optimization loops---a $+42.40\%$ absolute improvement over unregularized parametric linear baselines. Under stream-level mixed-task evaluations, MBH recovers GP-DR to $70.20\%$ accuracy ($+42.80\%$ recovery), showcasing a robust and theoretically sound foundation for dynamic modular deep learning.
\end{abstract}